\section{Информация о курсе} 
\subsection{Оценка} 
\[ 
\begin{aligned}
\text{О} = {} & 0.15 * \text{ДЗ} + \\
& 0.1 * \text{Quiz} + \\
& 0.15 * \text{Контест} + \\
& 0.2 * \text{КР (БЛОК)} + \\
& 0.25 * \text{Коллок (БЛОК)} + \\
& 0.05 * \text{Бонус от семинариста} + \\
& 0.2 * \text{Коллок по доп материалу} 
\end{aligned}
\]

\section{Введение}

\begin{definition}
    \textbf{Случайный эксперимент} -- это процесс, наблюдение или действие, результат которого невозможно точно предсказать заранее (до его проведения). Чтобы эксперимент можно было изучать методами теории вероятностей, он должен обладать следующими свойствами:

\begin{itemize}
    \item \textbf{Отсутствие детерминированной регулярности:} при многократном повторении эксперимента в неизменных условиях результаты могут различаться (исход единичного опыта случаен).
    \item \textbf{Воспроизводимость эксперимента:} эксперимент теоретически или практически можно повторить бесконечное число раз в одних и тех же контролируемых условиях. Его проведение и исход не должны быть завязаны на уникальные свойства или волю конкретной личности.
    \item \textbf{Статистическая устойчивость:} при очень большом числе повторений частота появления того или иного исхода стабилизируется и стремится к фиксированному числу (вероятности).
    \[
    \frac{N_{1, N}}{N_1} \approx \frac{N_{2, N}}{N_2} \approx \ldots \approx \frac{N_{k, N}}{N_k}
    \]
\end{itemize}
\end{definition}

Сопоставим каждому результату эксперимента свою $\omega $ --- элементарный исход

\[
    \{\omega \} = \Omega 
\]

$\Omega $ --- множество всех элементарных исходов

\begin{example}
    Кидаются 2 кости, нас интересует сумма:

\begin{enumerate}
    \item[1)] $\omega \in \{2; \dots; 12\},\ |\Omega| = 11$
    \item[2)] $\omega = (i_1, i_2),\ i_j \in \{1, \dots, 6\},\ |\Omega| = 36$
    \item[3)] $\omega = \{i_1, i_2\},\ |\Omega| = \frac{36-6}{2} + 6 = 21$
\end{enumerate}
\end{example}

\subsection{События}

\begin{definition}
    \textbf{Случайное событие} --- это подмножество пространства элементарных исходов: $\mathcal{A} \subset \Omega$.
\end{definition}

\begin{definition}
    \textbf{Вероятность} события $\mathcal{A}$ --- это предел относительной частоты его появления при стремлении числа экспериментов к бесконечности:
    \[
        \frac{N_{\mathcal{A}}}{N} \xrightarrow[N \to \infty]{} P(\mathcal{A})
    \]
\end{definition}

